\documentclass{article}
\usepackage[margin=25mm]{geometry}
\usepackage{amsmath,amssymb}
\usepackage{xurl}
\usepackage[colorlinks=true,linkcolor=blue,citecolor=blue,urlcolor=blue]{hyperref}
% Shared commands for notes in docs/paper-gaps/.

\DeclareUnicodeCharacter{2080}{\ifmmode{}_{0}\else\textsubscript{0}\fi}
\DeclareUnicodeCharacter{2081}{\ifmmode{}_{1}\else\textsubscript{1}\fi}
\DeclareUnicodeCharacter{2082}{\ifmmode{}_{2}\else\textsubscript{2}\fi}
\DeclareUnicodeCharacter{2099}{\ifmmode{}_{n}\else\textsubscript{n}\fi}

\newcommand{\Lean}{\textsc{Lean}}
\ExplSyntaxOn
\NewDocumentCommand{\leanid}{m}{
  \ifmmode
    \textsf{\detokenize{#1}}
  \else
    \group_begin:
    \urlstyle{sf}
    \tl_set:Nn \l_tmpa_tl {#1}
    \tl_replace_all:Nnn \l_tmpa_tl {\_} {_}
    \exp_args:NV \path \l_tmpa_tl
    \group_end:
  \fi
}
\ExplSyntaxOff
\providecommand{\tr}{\operatorname{tr}}
\newcommand{\tnleanIssueBase}{https://github.com/LionSR/TNLean/issues/}
\newcommand{\tnleanPrBase}{https://github.com/LionSR/TNLean/pull/}
\newcommand{\tnleanDocBase}{https://sirui-lu.com/TNLean/docs/}
\newcommand{\ghissue}[1]{\href{\tnleanIssueBase#1}{GitHub issue \##1}}
\newcommand{\ghpr}[1]{\href{\tnleanPrBase#1}{PR \##1}}
\newcommand{\doclink}[1]{\href{\tnleanDocBase#1.html}{\path{#1}}}

% Dirac notation and the identity operator, as in blueprint/src/macros/common.tex.
% \providecommand keeps note-local definitions authoritative.
\usepackage{dsfont}
\providecommand{\ket}[1]{|#1\rangle}
\providecommand{\bra}[1]{\langle #1|}
\providecommand{\braket}[2]{\langle #1 | #2 \rangle}
\providecommand{\ketbra}[2]{|#1\rangle\!\langle #2|}
\providecommand{\Id}{\mathds{1}}
\providecommand{\one}{\mathds{1}}
\providecommand{\C}{\mathbb{C}}
\providecommand{\MN}[1]{M_{#1}(\C)}
\providecommand{\MD}{M_D(\C)}

% Verdict marker: \gapnote{<kind>}{<status>}. Kind is the policy
% classification (clarification, local-correction, scope-restriction,
% unfaithful, false-source, open-gap); status is open, wip (elimination
% underway), resolved, or historical. Machine-read by the site index;
% prints a verdict line.
\newcommand{\gapnote}[2]{\par\noindent\textbf{Verdict:} #1 (#2).\par}

\providecommand{\gapnote}[2]{\par\noindent\textbf{Verdict:} #1 (#2).\par}
\title{Schmidt Rank of the ALVV Frustration-Free Truncation}
\author{}
\date{October 29, 2026}
\begin{document}
\maketitle
\gapnote{false-source}{open}

\section{The rank assertion}

Arad, Landau, Vazirani and Vidick~\cite{Arad2016RigorousRG},
\href{https://arxiv.org/abs/1602.08828v2}{arXiv:1602.08828v2},
Theorem \texttt{thm:FF-tH}, item 2, asserts that the frustration-free
truncated Hamiltonian has operator Schmidt rank at most $d^2+2$ across
every cut. The definition and argument are in
\texttt{local\_alg\_arxiv\_v2.tex}, lines 1467--1493 and 1543--1568.
This constant is false for an interval with nonempty complementary
parts on both sides. The example below has $d=2$ and rank seven,
with unique ground state and spectral gap exactly one both before
and after truncation. The corrected general bound is $d^2+3$.

Throughout, the terms truncated on an interval are its internal bonds,
as in the explicit sum defining $H_J$ in Definition \texttt{def:FF-tH}.
The subsequent parity-index notation in that definition includes an
extra endpoint; that separate indexing typo is not used here.
Operator Schmidt rank means the least number of products $A_j\otimes B_j$
whose sum is the operator. This note gives an exact mathematical
counterexample, not a formalized Lean counterexample theorem.

\section{An eight-qubit chain}

Number the sites $0,\ldots,7$ and let $p=|0\rangle\langle0|$.
For bonds $i=0,\ldots,6$, set
\[
 h_i=I-p_i p_{i+1}\quad(i\ne3),\qquad
 h_3=|\psi^-\rangle\langle\psi^-|,\qquad
 |\psi^-\rangle=(|01\rangle-|10\rangle)/\sqrt2.
\]
Each term is a nonzero orthogonal projection of norm one, supported on
its two assigned sites. The four even bonds are disjoint and cover
all sites. Thus
\[
 H=\sum_{i=0}^6h_i\ \ge\ h_0+h_2+h_4+h_6
       \ \ge\ I-|0^8\rangle\langle0^8|.
\]
All seven terms annihilate $|0^8\rangle$. The vector $|10000000\rangle$
is an eigenvector of energy one. Hence $H$ is frustration free, its
ground state is unique, and its spectral gap is exactly one.

Take the region $J=\{1,\ldots,6\}$, with internal bonds $1,\ldots,5$.
The source's truncated Hamiltonian is
\[
 \widetilde H=2I-Q_o-Q_e+h_0+h_6,
 \qquad
 Q_o=(I-h_1)(I-h_3)(I-h_5),\quad
 Q_e=(I-h_2)(I-h_4).
\]
The terms in each product have disjoint supports.
Furthermore, $h_0$, $I-Q_e$, and $h_6$ are projections on disjoint
blocks covering the eight sites, with common kernel spanned by $|0^8\rangle$.
Their sum dominates $I-|0^8\rangle\langle0^8|$.
Since $I-Q_o\ge0$ and $|10000000\rangle$ still has energy one,
$\widetilde H$ also has unique ground state $|0^8\rangle$ and gap one.

\section{The exact Schmidt rank}

Cut between sites 3 and 4. On the left and right four-qubit factors put
\begin{align*}
 E_L&=I\otimes p\otimes p\otimes I,&
 E_R&=I\otimes p\otimes p\otimes I,\\
 O_L&=I\otimes I\otimes p\otimes p,&
 O_R&=p\otimes p\otimes I\otimes I,\\
 B_L&=I-p\otimes p\otimes I\otimes I,&
 B_R&=I-I\otimes I\otimes p\otimes p.
\end{align*}
Here the identity in an additive expression acts on all four sites.
Let
\[
 X=\begin{pmatrix}0&1\\1&0\end{pmatrix},\quad
 W=\begin{pmatrix}0&1\\-1&0\end{pmatrix},\quad
 Z=\begin{pmatrix}1&0\\0&-1\end{pmatrix}.
\]
For $\sigma\in\{I,X,W,Z\}$, define
$E_{L,\sigma}=I\otimes p\otimes p\otimes\sigma$ and
$E_{R,\sigma}=\sigma\otimes p\otimes p\otimes I$.
The identity
\[
 I-h_3=\tfrac14(3I\otimes I+X\otimes X-W\otimes W+Z\otimes Z)
\]
gives
\begin{align}
 \widetilde H={}&-\tfrac34E_{L,I}\otimes E_{R,I}
 -\tfrac14E_{L,X}\otimes E_{R,X}
 +\tfrac14E_{L,W}\otimes E_{R,W}
 -\tfrac14E_{L,Z}\otimes E_{R,Z}\notag\\
 &-O_L\otimes O_R+(2I+B_L)\otimes I+I\otimes B_R.
 \label{eq:ff_rank_seven}
\end{align}
Both ordered seven-element families in this expression are linearly
independent over $\mathbb C$. To check this directly, number the
four-qubit computational basis $0,\ldots,15$. In the left family select
matrix entries
\[
 (0,0),(0,1),(1,0),(1,1),(4,4),(2,2),(5,5);
\]
in the right family select
\[
 (0,0),(0,8),(8,0),(8,8),(2,2),(4,4),(1,1).
\]
The resulting matrices, whose columns follow the order in
\eqref{eq:ff_rank_seven}, are
\[
 \begin{pmatrix}
 1&0&0&1&1&2&1\\0&1&1&0&0&0&0\\0&1&-1&0&0&0&0\\
 1&0&0&-1&0&2&1\\0&0&0&0&1&3&1\\0&0&0&0&0&2&1\\0&0&0&0&0&3&1
 \end{pmatrix},\qquad
 \begin{pmatrix}
 1&0&0&1&1&1&0\\0&1&1&0&0&0&0\\0&1&-1&0&0&0&0\\
 1&0&0&-1&0&1&0\\0&0&0&0&1&1&1\\0&0&0&0&0&1&0\\1&0&0&1&1&1&1
 \end{pmatrix}.
\]
Their determinants are $-4$ and $4$. All seven coefficients in
\eqref{eq:ff_rank_seven} are nonzero, so the operator Schmidt rank
is exactly seven. This exceeds the printed bound $d^2+2=6$.

\section{The corrected count and the formal boundary}

For a cut in the truncated region, exactly one parity product can contain
a bond crossing the cut. Its operator Schmidt rank is at most $d^2$:
all its other factors act entirely on one side. The other parity product
has rank at most one. Finally,
\[
 2I+H_{\overline J}= (2I+H_{\mathrm{left}})\otimes I
                         +I\otimes H_{\mathrm{right}}
\]
has rank at most two. Thus $\operatorname{OSR}(\widetilde H)\le d^2+3$.
The example saturates this count. The source's interior-cut argument
counts the complementary Hamiltonian as one contribution, whereas it can
have independent terms on both sides. At cuts outside the region the
usual $d^2+2$ bound still holds.

The correction changes an additive constant in the rank estimate; it
does not invalidate the ground-space or spectral-gap conclusions of the
frustration-free truncation theorem. The physical detectability and
range-projection reductions are separately proved. The regional truncated
Hamiltonian and its corrected rank estimate remain to be formalized under
\href{https://github.com/LionSR/TNLean/issues/9833}{the ALVV programme}.
No source-labelled formalization of the printed rank assertion is claimed.

\bibliographystyle{unsrt}
\bibliography{references}
\end{document}
